% Fibonacci-Pell nearest gaps: an all-exponent classification and
% quadratic-unit orbit rigidity
%
% Authors:
%   Dr. Denys Dutykh (Mathematics Department, Khalifa University of Science
%   and Technology, Abu Dhabi, UAE)
%   Prof. Laurent Vuillon (Univ. Savoie Mont Blanc, CNRS, LAMA, Chambery,
%   France)
%
\section{Leading unit exponents in the even branch}
\label{sec:leading}

Assume throughout \cref{sec:leading,sec:height,sec:classification} that an
even-exponent target is fixed, and retain $j$ and $M=n-j$ from
\cref{lem:normalisation}. By \cref{prop:parity-selector}, $3\nmid q$ and
$M$ is odd; in particular, \cref{lem:rectangle} applies. The first task is
to show that every target in the unresolved range has a genuinely negative
normalised exponent.

\begin{lemma}[Ratio bounds]
\label{lem:ratio-bounds}
For $q\geqslant3$,

\begin{equation}
 \lambda B_q<A_q<3B_q.
 \label{eq:AB-ratio}
\end{equation}
\end{lemma}

\begin{proof}
Using $F_{q+1}=F_q+F_{q-1}$, direct expansion gives

\begin{equation}
 A_q-\lambda B_q=\sqrt2F_q-2F_{q-1}>0,
\end{equation}

because $F_q/F_{q-1}\geqslant3/2>\sqrt2$ for $q\geqslant3$.
Similarly,

\begin{equation}
 3B_q-A_q=2\bigl(\sqrt2F_{q+1}-2F_q\bigr)>0.
\end{equation}

Indeed, $F_{q+1}/F_q\geqslant3/2>\sqrt2$ for $q\geqslant3$, which proves
the last inequality.
\end{proof}

\begin{proposition}[Negative-exponent reduction]
\label{prop:negative-exponent}
Every even-exponent target with $q\geqslant7$ has

\begin{equation}
 j=-h,
 \qquad
 h\geqslant3\ \text{odd}.
 \label{eq:j-negative}
\end{equation}
\end{proposition}

\begin{proof}
On the plus branch, \cref{eq:plus-orbit} has positive rational coefficient
$W=C_M-g$. Quadratic conjugation gives

\begin{equation}
 2W=A_q^2\lambda^{-j}-B_q^2\lambda^j>0.
\end{equation}

Thus $\lambda^j<A_q/B_q<3<\lambda^2$. Since $j$ is odd, $j\leqslant1$.
On the minus branch the actual rational coefficient is
$C_M-g=-W>0$. Conjugating \cref{eq:minus-orbit} gives

\begin{equation}
 2(C_M-g)=B_q^2\lambda^{-j}-A_q^2\lambda^j>0,
\end{equation}

and hence

\begin{equation}
 \lambda^j<\frac{B_q}{A_q}<\lambda^{-1}.
\end{equation}

Hence the odd integer $j$ is at most $-3$ there.

It remains to remove $j=1$ and $j=-1$ on the plus branch. If $j=1$,
comparison of the $\sqrt2$-coefficients in \cref{eq:plus-orbit} gives

\begin{equation}
 P_M=F_{q-1}^2+F_{q+1}^2
 =3F_q^2+2(-1)^q.
 \label{eq:j1-recurrence}
\end{equation}

This recurrence equation has exactly

\begin{equation}
 (q,M)=(1,1),(2,3),(4,5)
 \label{eq:j1-solutions}
\end{equation}

among nonnegative $q$ and positive odd $M$. This follows directly from the
complete near-square classification of Alekseyev and Tengely
\cite[Table~1]{AlekseyevTengely2014}. In their notation the Pell sequence is
$U_M(2,-1)$. They prove that all Pell values of the forms $3y^2+2$ and
$3y^2-2$ are, respectively,

\begin{equation}
\{2,5,29\}
\qquad\text{and}\qquad
\{1\}.
\label{eq:near-square-values}
\end{equation}

If $q$ is even, \cref{eq:j1-recurrence} has the first form in
\cref{eq:near-square-values}. The value $2=P_2$ is removed because $M$ is
odd, while $5=P_3$ and $29=P_5$ give $F_q=1$ and $F_q=3$. Even parity of
$q$ leaves $(q,M)=(2,3),(4,5)$. If $q$ is odd, the second form in
\cref{eq:near-square-values} gives $P_M=1$, hence $M=1$ and $F_q=1$;
odd parity leaves $(q,M)=(1,1)$. Direct substitution proves the converse.

The nearest plus window is absent at $q=2$, and $q=4$ is the low plus target
in \cref{eq:two-solutions}. Thus no $j=1$ target has $q\geqslant7$.

If $j=-1$, the radical coefficient instead becomes

\begin{equation}
 P_M=F_q^2+2F_{q+1}^2+2F_qF_{q+1}
 =F_{q+1}^2+F_{q+2}^2
 =F_{2q+3}.
 \label{eq:jminus1}
\end{equation}

Alekseyev's complete Fibonacci--Pell intersection theorem states that the
common values of the two sequences are $0,1,2,5$ \cite[Theorem~9]{Alekseyev2011}.
Since $2q+3\geqslant17$ for $q\geqslant7$, \cref{eq:jminus1} is
impossible. All remaining exponents therefore have the form~\labelcref{eq:j-negative}.
\end{proof}

\begin{remark}[Independent verification of the near-square step]
\label{rem:quartic-status}
The published theorem \cite{AlekseyevTengely2014} is the completeness input;
this boundary classification is not new here. As an independent audit, the
negative-Pell identity converts \cref{eq:j1-recurrence}, with
$y=F_q$ and $s=(-1)^q$, into

\begin{equation}
 C_M^2=18y^4+24sy^2+7.
\end{equation}

After $y=z+1$, two complete Magma V2.29-9 integral-point computations give
exactly the points corresponding to \cref{eq:j1-solutions}. The supplement
contains both inputs, complete transcripts, and independent arithmetic
verifiers. These computations corroborate the published theorem and are not
used as an unrecorded bounded search.
\end{remark}

Thus every even-exponent target in the unresolved range has primitive unit
direction $\lambda^{-h}$ with $h\geqslant3$ odd. The qualitative
leading-unit ambiguity has disappeared; the next task is quantitative,
namely to separate the negative exponent $h$ from the nearest Pell exponent
$n$.
